% ======================================================================================
% Nom     : reports/latex/sections/08_analyse_critique.tex
% Rôle    : Section d'analyse critique des résultats.
% Auteur  : Maxime BRONNY
% Version : V1
% Cadre   : Réalisé dans le cadre du cours Techniques d'apprentissage artificiel
%           Master 1 Informatique Big Data
% Usage   : Fichier inclus dans main.tex lors de la compilation du rapport.
% ======================================================================================

\section{Analyse critique}
\label{sec:analyse}

\subsection{Pourquoi l'arbre de décision domine-t-il ?}

La supériorité de l'arbre s'explique par la nature du problème. Le risque de
crédit repose sur des \emph{interactions} entre variables : un taux
d'endettement de $30\,\%$ n'a pas la même signification pour un dossier de
grade A que pour un grade F, et un antécédent de défaut ne pèse pas pareil
selon le revenu. La régression logistique, qui additionne des contributions
indépendantes des variables, ne peut pas représenter ces effets croisés ; sa
frontière de décision reste un hyperplan, et son AUC plafonne à $0{,}848$.
L'arbre, en enchaînant des tests conditionnels (d'abord le grade, puis le taux
d'endettement \emph{dans} chaque branche, etc.), modélise naturellement ces
interactions, d'où son AUC de $0{,}909$ et surtout sa précision de
$0{,}977$. Le $k$-NN, qui capte la structure locale des données sans la
modéliser explicitement, se place logiquement entre les deux.

\subsection{L'apport mesurable de l'ensemble de validation}

Le réglage du seuil de décision sur la validation illustre concrètement
l'intérêt du découpage en trois ensembles. Pour la régression logistique, le
seuil retenu ($0{,}3$) relève le rappel à $0{,}69$ là où le seuil naïf de
$0{,}5$ (celui de la première version du projet) donnait un F1 d'environ
$0{,}50$ ; le F1 atteint désormais $0{,}626$, soit un gain d'environ douze
points obtenu \emph{sans toucher au modèle}, simplement par un protocole
expérimental plus rigoureux. À l'inverse, l'arbre préfère un seuil élevé
($0{,}6$) qui filtre les feuilles incertaines et lui donne sa précision
remarquable. Le seuil optimal est donc bien une propriété propre à chaque
modèle, qu'il serait incorrect de fixer uniformément à $0{,}5$.

\subsection{Surapprentissage : un risque contrôlé}

La comparaison des métriques d'entraînement et de test
(\texttt{results/metrics.json}) montre des écarts très faibles : F1 de
$0{,}623$ (train) contre $0{,}626$ (test) pour la régression logistique,
$0{,}814$ contre $0{,}817$ pour l'arbre. Deux mécanismes y contribuent : les
garde-fous de l'arbre (taille minimale des n\oe uds et des feuilles) et surtout
la sélection de la profondeur sur la validation, qui a écarté la profondeur 9
(F1 de validation $0{,}817$ contre $0{,}819$ pour la profondeur 7) ; le début
de surapprentissage a donc été détecté et évité par le protocole lui-même. Le
$k$-NN présente l'écart le plus visible (AUC $0{,}907$ en train contre
$0{,}877$ en test), comportement attendu d'une méthode à base d'instances
évaluée sur ses propres exemples mémorisés ; l'écart reste modéré grâce au
lissage d'un $k$ élevé ($31$).

\subsection{Le rappel, limite commune aux trois modèles}

Le résultat le plus instructif n'est pas la hiérarchie des modèles mais leur
plafond commun : environ $30\,\%$ des défauts ne sont détectés par aucun
modèle, alors même que leurs familles algorithmiques sont très différentes.
Cela suggère que la limite ne vient pas des algorithmes mais de
l'\emph{information disponible} : une partie des défauts n'est simplement pas
prévisible à partir des onze variables du dossier de crédit (accidents de la
vie, perte d'emploi postérieure à l'octroi, etc.). Pour une banque, ce constat
est important : améliorer encore le modèle rapporterait moins que collecter des
variables plus informatives.

\subsection{Points forts et réserves sur la qualité des données}

Au crédit du travail : un protocole sans fuite de données (paramètres de
prétraitement appris sur le seul entraînement), une évaluation finale unique,
des implémentations validées par une référence externe (écart d'AUC
$\leq 5 \times 10^{-4}$), une expérience entièrement reproductible (graine
fixe) et couverte par onze tests automatisés.

Côté réserves : l'imputation par la médiane de \texttt{loan\_int\_rate}
(près de $10\,\%$ de valeurs manquantes) lisse l'information d'une variable
probablement liée au risque ; le label encoding impose un ordre artificiel aux
deux variables nominales, ce qui pénalise principalement la régression
logistique, et une partie de son retard pourrait venir de là plutôt que de sa
linéarité seule ; enfin, le déséquilibre des classes n'est traité que par le
réglage du seuil, sans ré-échantillonnage ni pondération.
